% =============================================================
% บทที่ 2 - ทฤษฎีและงานวิจัยที่เกี่ยวข้อง
% =============================================================

% -----------------------------------------------------------
\section{ทฤษฎีที่เกี่ยวข้อง}
% -----------------------------------------------------------

\subsection{เว็บสแครปปิง (Web Scraping)}

เว็บสแครปปิงคือกระบวนการดึงข้อมูลจากเว็บไซต์โดยอัตโนมัติ แทนการคัดลอกข้อมูลด้วยมือ โดยทั่วไปจะส่งคำขอ HTTP ไปยังเว็บเซิร์ฟเวอร์แล้ววิเคราะห์โครงสร้าง HTML ของหน้าเว็บเพื่อสกัดข้อมูลที่ต้องการ~\cite{trafilatura2021} เครื่องมือที่นิยมใช้ ได้แก่ ไลบรารี BeautifulSoup สำหรับวิเคราะห์และค้นหาองค์ประกอบในเอกสาร HTML ด้วย CSS Selector และไลบรารี httpx สำหรับส่งคำขอ HTTP แบบอะซิงโครนัส อย่างไรก็ตาม เว็บไซต์บางแห่งแสดงผลเนื้อหาผ่าน JavaScript ต้องใช้เครื่องมือจำลองเบราว์เซอร์ เช่น Playwright เพื่อให้ได้ข้อมูลที่แสดงผลสมบูรณ์

นอกจากนี้ เว็บไซต์หลายแห่งยังเผยแพร่ข่าวผ่าน RSS Feed ซึ่งเป็นรูปแบบ XML มาตรฐานที่ใช้แจกจ่ายเนื้อหาที่อัปเดตเป็นประจำ ทำให้สามารถดึงหัวข้อข่าว ลิงก์ และบทสรุปสั้น ๆ ได้โดยไม่ต้องวิเคราะห์โครงสร้าง HTML โดยตรง~\cite{rss2009} ในระบบนี้ใช้ทั้งสองเทคนิคร่วมกันตามความเหมาะสมของแต่ละสำนักข่าว

\subsection{การสรุปข้อความด้วยโมเดลภาษาขนาดใหญ่}

การสรุปข้อความ (Text Summarization) คือกระบวนการย่อข้อความต้นฉบับให้สั้นลงโดยยังคงประเด็นสำคัญไว้ โดยแบ่งเป็น 2 แนวทางหลัก ได้แก่ การสรุปแบบ Extract (เลือกประโยคสำคัญจากต้นฉบับ) และการสรุปแบบ Abstractive (เขียนข้อความสรุปขึ้นใหม่ด้วยคำของตนเอง)~\cite{allahyari2017} โมเดลภาษาขนาดใหญ่ (Large Language Model: LLM) เช่น GPT และ Gemini ได้รับการฝึกบนคลังข้อมูลจำนวนมหาศาล จึงมีความสามารถในการสรุปแบบ Abstractive ที่เป็นธรรมชาติและแม่นยำ~\cite{gemini2024}

ในระบบนี้ใช้โมเดลภาษาขนาดใหญ่ผ่านบริการ KKU Gen.AI โดยกำหนด System Prompt ให้โมเดลอ่านเนื้อหาข่าวในรูปแบบ Markdown แล้วส่งคืนบทสรุป ประเด็นสำคัญ อารมณ์ของข่าว และคำสำคัญ ในรูปแบบ JSON ที่มีโครงสร้างแน่นอน ทำให้ระบบสามารถแปลงผลลัพธ์ไปเก็บในฐานข้อมูลได้โดยตรง~\cite{openai2024}

\subsection{การจำแนกข้อความภาษาไทย}

การจำแนกข้อความ (Text Classification) เป็นงานประมวลผลภาษาธรรมชาติที่จัดกลุ่มข้อความตามเนื้อหา ภาษาไทยเป็นภาษาที่ไม่มีการเว้นวรรคระหว่างคำ จึงจำเป็นต้องใช้เครื่องมือตัดคำภาษาไทย เช่น PyThaiNLP ซึ่งเป็นไลบรารีประมวลผลภาษาธรรมชาติภาษาไทย~\cite{pythainlp2021}

ระบบใช้แนวทางไฮบริดระหว่างกฎคำสำคัญ (Keyword-based Rules) กับการเรียนรู้ของเครื่อง โดยกฎคำสำคัญให้คะแนนตามการพบคำที่เกี่ยวข้องในแต่ละหมวดหมู่ พร้อมน้ำหนักของคำที่ยาวและกฎเชิงบริบทร่วม (Compound Rules) ในขณะที่โมเดล SVM ที่ได้รับการฝึกด้วยชุดข้อมูลข้อความภาษาไทยโดยใช้ TF-IDF ทำหน้าที่ตรวจสอบบริบทของข้อความ และระบบจะตัดสินใจใช้ผลจากวิธีที่เหมาะสมที่สุด~\cite{scikit2024}

\subsection{การจัดการแคชและคุกกี้}

แคช (Cache) คือกลไกการจัดเก็บสำเนาข้อมูลที่ได้จากการดึงหรือคำนวณไว้ชั่วคราว เพื่อให้การเรียกใช้ครั้งถัดไปสามารถตอบสนองได้รวดเร็วขึ้นโดยไม่ต้องทำงานซ้ำ สำหรับระบบรวบรวมข่าวอัตโนมัติที่ทำงานเป็นรอบ การใช้แคชช่วยให้ระบบไม่ต้องดึงหน้าเว็บและบทความเดิมซ้ำในแต่ละรอบ โดยตรวจสอบจากตัวระบุเฉพาะของข่าว เช่น URL ก่อนตัดสินใจว่าจะต้องดึงเนื้อหาใหม่หรือใช้ข้อมูลเดิมจากแคช ทำให้ลดภาระของเว็บไซต์ต้นทางและประหยัดทรัพยากรของเซิร์ฟเวอร์

คุกกี้ (Cookie) คือข้อมูลขนาดเล็กที่เว็บเซิร์ฟเวอร์ส่งให้เบราว์เซอร์เพื่อใช้จดจำสถานะของผู้ใช้หรือเซสชันการเชื่อมต่อ เว็บไซต์ข่าวบางแห่งกำหนดให้เบราว์เซอร์ต้องมีคุกกี้ที่ถูกต้องหรือผ่านการตรวจสอบ เช่น หน้าเพจตรวจสอบของ Cloudflare ก่อนจึงจะแสดงเนื้อหาให้ ระบบนี้จึงต้องจัดการคุกกี้ของเบราว์เซอร์ผ่านการจำลองเบราว์เซอร์ด้วย Playwright~\cite{playwright2024} และจัดเก็บคุกกี้ไว้ใช้ในการดึงข้อมูลครั้งถัดไป เพื่อให้สามารถดึงข่าวจากเว็บไซต์ดังกล่าวได้อย่างถูกต้องและต่อเนื่องโดยไม่ต้องเริ่มต้นเซสชันใหม่ทุกครั้ง

นอกจากการจัดการคุกกี้เพื่อการดึงข้อมูลแล้ว ระบบยังใช้คุกกี้เพื่อวัตถุประสงค์ในการปรับแต่งประสบการณ์ของผู้ใช้อีกสองประการ ประการแรกคือ \textbf{คุกกี้เพื่อการปรับแต่งเฉพาะบุคคล (Personalize Cookies)} ซึ่งใช้จัดเก็บค่ากำหนดของผู้ใช้ เช่น หมวดหมู่ข่าวที่ต้องการติดตามเป็นพิเศษ ภาษาในการแสดงผล รูปแบบการแสดงข่าว หรือรายการข่าวที่ผู้ใช้บันทึกไว้ เพื่อให้ผู้ใช้แต่ละคนได้รับมุมมองและเนื้อหาที่ตรงกับความต้องการของตนเองเมื่อเข้าใช้งานระบบในครั้งถัดไป ประการที่สองคือ \textbf{คุกกี้เพื่ออัลกอริทึม (Algorithm Cookies)} ซึ่งใช้บันทึกพฤติกรรมการอ่านข่าว เช่น หมวดหมู่ข่าวที่ผู้ใช้คลิกอ่านบ่อย ข่าวที่ใช้เวลาอ่านนาน หรือข่าวที่ผู้ใช้กดบันทึกหรือแชร์ ข้อมูลพฤติกรรมเหล่านี้จะถูกส่งให้อัลกอริทึมแนะนำข่าว (Recommendation Algorithm) เพื่อวิเคราะห์ความสนใจและคัดสรรข่าวที่เกี่ยวข้องมากที่สุดมาแสดงเป็นลำดับแรกบนหน้าเว็บ ซึ่งช่วยเพิ่มความเกี่ยวข้องของเนื้อหาและประสบการณ์การใช้งานที่ดียิ่งขึ้นสำหรับผู้ใช้แต่ละคน

\subsection{อัลกอริทึมการกำหนดคุกกี้ (Assign Cookies Algorithm)}

เพื่อให้คุกกี้สามารถสะท้อนความสนใจของผู้ใช้ได้อย่างเป็นปัจจุบัน ระบบจึงออกแบบอัลกอริทึมการกำหนดคุกกี้ที่ทำงานเป็นลำดับขั้นตอน ดังนี้

\begin{enumerate}[label=\arabic*)]
    \item \textbf{สร้างคุกกี้เริ่มต้น (Initialization):} เมื่อผู้ใช้เข้าใช้งานครั้งแรก ระบบสร้างตัวระบุเฉพาะของผู้ใช้ (Anonymous Cookie ID) และกำหนดคุกกี้เริ่มต้น โดยค่าเริ่มต้นของความสนใจแต่ละหมวดหมู่ถูกตั้งไว้เท่ากัน เพื่อให้ผลการแนะนำไม่เอนเอียง
    \item \textbf{บันทึกพฤติกรรมการอ่าน (Logging):} ทุกครั้งที่ผู้ใช้คลิกเปิดข่าว ใช้เวลาในการอ่าน กดบันทึก หรือแชร์ข่าว ระบบจะบันทึกเหตุการณ์ (Event) พร้อมข้อมูลหมวดหมู่ของข่าวนั้น และเพิ่มความถี่ที่สะสมของแต่ละหมวดหมู่ลงในคุกกี้เพื่ออัลกอริทึม
    \item \textbf{คำนวณคะแนนความสนใจ (Interest Score):} อัลกอริทึมคำนวณคะแนนความสนใจของแต่ละหมวดหมู่โดยถ่วงน้ำหนักจากความถี่ที่เข้าชม เวลารวมที่ใช้ในการอ่าน และการกระทำที่แสดงความสนใจสูง เช่น การกดบันทึกหรือแชร์ โดยหมวดที่ถูกอ่านถี่และใช้เวลานานจะได้รับคะแนนสูง
    \item \textbf{กำหนดคุกกี้ปรับแต่งเฉพาะบุคคล (Assign Personalize Cookie):} เมื่อคะแนนความสนใจของหมวดหมู่ใดถึงเกณฑ์ที่กำหนด อัลกอริทึมจะอัปเดตคุกกี้ปรับแต่งให้รวมหมวดหมู่ดังกล่าวไว้ในรายการหมวดที่ผู้ใช้สนใจเป็นพิเศษ และปรับลำดับการนำเสนอข่าวและภาษาที่เหมาะสมให้สอดคล้องกับค่ากำหนดของผู้ใช้
    \item \textbf{กำหนดคุกกี้อัลกอริทึมแนะนำข่าว (Assign Algorithm Cookie):} ระบบนำคะแนนความสนใจที่คำนวณได้ทั้งหมดไปเข้ารหัสและจัดเก็บลงในคุกกี้เพื่ออัลกอริทึม โดยคุกกี้นี้ทำหน้าที่เป็นข้อมูลนำเข้าของบริการแนะนำข่าว (Recommendation) ใช้เรียงลำดับข่าวจากหมวดที่มีความสนใจสูงไปต่ำก่อนแสดงบนหน้าเว็บ และใช้เป็นน้ำหนักในการจัดอันดับข่าวแบบเรียลไทม์ในครั้งถัดไป
    \item \textbf{ป้องกันการเอนเอียงของข้อมูล (防止 Filter Bubble):} เพื่อไม่ให้ผู้ใช้เห็นเฉพาะข่าวที่ระบบคิดว่าชอบจนเกิดปรากฏการณ์ Filter Bubble อัลกอริทึมจะนำเสนอข่าวจากหมวดหมู่อื่นในสัดส่วนที่กำหนดไว้ด้วย พร้อมทั้งกำหนดอายุการหมดอายุของคุกกี้ให้สั้น (Session Expiry) เพื่อไม่ให้คุกกี้คงอยู่ตลอดไปและลดความเสี่ยงด้านความเป็นส่วนตัว
\end{enumerate}

การทำงานของอัลกที่ริทึมการกำหนดคุกกี้ข้างต้นเป็นการผสานระหว่างคุกกี้เพื่อการปรับแต่งเฉพาะบุคคล คุกกี้เพื่ออัลกอริทึม และบริการแนะนำข่าวเข้าไว้ด้วยกัน เพื่อให้ระบบสามารถปรับเปลี่ยนประสบการณ์การอ่านข่าวให้ตรงกับความสนใจของผู้ใช้แต่ละคนได้อย่างต่อเนื่องและปลอดภัย

% -----------------------------------------------------------
\section{งานวิจัยที่เกี่ยวข้อง}
% -----------------------------------------------------------

ในส่วนนี้จะนำเสนองานวิจัยหรือระบบที่มีลักษณะคล้ายคลึงกับโครงงานนี้ เพื่อเป็นแนวทางในการพัฒนา โดยพิจารณาครอบคลุมในด้านการสกัดเนื้อหาจากเว็บ การสรุปข้อความ และการจำแนกหมวดหมู่ข่าว

\subsection{งานสกัดเนื้อหาบทความจากเว็บ}

Barbaresi \cite{trafilatura2021} ได้พัฒนาไลบรารี Trafilatura สำหรับสกัดข้อความหลักออกจากหน้าเว็บโดยอัตโนมัติ โดยไม่ขึ้นกับเว็บไซต์ใดเว็บไซต์หนึ่ง ไลบรารีนี้รองรับการจัดรูปแบบ HTML ที่หลากหลายและคืนข้อความในรูปแบบ Markdown ซึ่งเป็นประโยชน์โดยตรงต่อการนำเนื้อหาไปใช้กับโมเดลภาษาขนาดใหญ่ อย่างไรก็ตาม วิธีนี้มีข้อจำกัดเมื่อเนื้อหาถูกเรนเดอร์ด้วย JavaScript บางกรณี ทำให้ระบบนี้ต้องเสริมด้วยการจำลองเบราว์เซอร์ด้วย Playwright สำหรับเว็บไซต์กลุ่มดังกล่าว

\subsection{งานสรุปข้อความอัตโนมัติ}

Allahyari และคณะ \cite{allahyari2017} ได้สำรวจเทคนิคการสรุปข้อความอย่างเป็นระบบทั้งแบบ Extract และ Abstractive และชี้ให้เห็นว่าการสรุปแบบ Abstractive ที่ใช้การเรียนรู้เชิงลึกให้ผลลัพธ์ที่เป็นธรรมชาติกว่าอย่างมีนัยสำคัญ แต่ต้องใช้ทรัพยากรในการคำนวณสูง ด้วยความก้าวหน้าของโมเดลภาษาขนาดใหญ่ \cite{gemini2024} การสรุปแบบ Abstractive จึงเข้าถึงได้ง่ายผ่าน API โดยไม่ต้องฝึกโมเดลเอง โครงงานนี้จึงเลือกใช้แนวทางดังกล่าวเพื่อสรุปข่าวเป็นภาษาเดียวกับต้นฉบับโดยไม่ต้องใช้ทรัพยากรฝึกโมเดลของตนเอง

\subsection{งานจำแนกหมวดหมู่ข่าวภาษาไทย}

งานวิจัยด้านการจำแนกข้อความภาษาไทยมักประสบปัญหาการเว้นวรรคคำ จึงนิยมใช้เครื่องมือตัดคำภาษาไทย เช่น PyThaiNLP \cite{pythainlp2021} ร่วมกับโมเดลแมชชีนเลิร์นนิง เช่น SVM กับคุณลักษณะ TF-IDF ซึ่งเป็นแนวทางที่มีประสิทธิภาพและใช้ทรัพยากรต่ำ โครงงานนี้ประยุกต์ใช้แนวทางไฮบริด โดยใช้กฎคำสำคัญที่มีน้ำหนักและลำดับความสำคัญของ URL สำหรับการจำแนกแบบรวดเร็ว และใช้โมเดล SVM เป็นตัวตรวจสอบบริบทสำหรับหมวดหมู่ที่กฎคำสำคัญมีความแม่นยำไม่เพียงพอ

\subsection{สรุปผลการเปรียบเทียบงานวิจัยที่เกี่ยวข้อง}

จากการทบทวนวรรณกรรมข้างต้น สามารถสรุปผลการเปรียบเทียบคุณสมบัติได้ดังตารางที่ \ref{tab:compare_research}

\begin{table}[H]
\centering
\caption{เปรียบเทียบงานวิจัยที่เกี่ยวข้อง}
\label{tab:compare_research}
\begin{tabular}{|l|c|c|c|c|c|c|}
\hline
\textbf{งานวิจัย} & \textbf{สกัดเนื้อหา} & \textbf{สรุป LLM} & \textbf{จำแนกหมวด} & \textbf{Realtime} & \textbf{แคช/คุกกี้} & \textbf{อัลกอริทึมแนะนำ} \\
\hline
Trafilatura \cite{trafilatura2021} & \cmark & -- & -- & -- & -- & -- \\ \hline
งานสรุปข้อความ \cite{allahyari2017} & -- & \cmark & -- & -- & -- & -- \\ \hline
งานจำแนกข้อความภาษาไทย \cite{pythainlp2021} & -- & -- & \cmark & -- & -- & -- \\ \hline
\textbf{โครงงานนี้} & \cmark & \cmark & \cmark & \cmark & \cmark & \cmark \\
\hline
\end{tabular}
\end{table}

จากตารางที่ \ref{tab:compare_research} พบว่ายังไม่มีงานใดที่ผสมผสานการสกัดเนื้อหา การสรุปด้วยโมเดลภาษาขนาดใหญ่ การจำแนกหมวดหมู่ การแสดงผลแบบเรียลไทม์ การจัดการแคชและคุกกี้ และอัลกอริทึมแนะนำข่าวในระบบเดียวครบทุกด้าน โครงงานนี้จึงมุ่งพัฒนาเว็บแอปพลิเคชันที่รวบรวมคุณสมบัติทั้งหมดไว้ในระบบเดียวอย่างครบถ้วน

\subsection{เทคโนโลยีและเครื่องมือที่เกี่ยวข้อง}

\subsubsection{FastAPI}

FastAPI เป็นเฟรมเวิร์คภาษา Python สำหรับพัฒนา REST API ที่มีประสิทธิภาพสูง รองรับการเขียนแบบอะซิงโครนัส (Asynchronous) และสร้างเอกสาร API อัตโนมัติ ใช้สำหรับสร้าง REST API และให้บริการไฟล์ด้านหน้าเว็บในระบบนี้

\subsubsection{Playwright}

Playwright เป็นเครื่องมือสำหรับควบคุมเบราว์เซอร์ Chromium ผ่านโปรแกรม ใช้สำหรับดึงข้อมูลจากเว็บไซต์ที่แสดงผลเนื้อหาผ่าน JavaScript ซึ่งวิธี HTTP ทั่วไปไม่สามารถเข้าถึงได้

\subsubsection{PyThaiNLP}

PyThaiNLP เป็นไลบรารีประมวลผลภาษาธรรมชาติภาษาไทย ประกอบด้วยเครื่องมือตัดคำ แปลงตัวเลข และข้อมูลพจนานุกรมภาษาไทย ใช้ในการตัดคำภาษาไทยเพื่อการจำแนกหมวดหมู่ข่าว

\subsubsection{Socket.IO}

Socket.IO เป็นไลบรารีสำหรับการสื่อสารแบบเรียลไทม์แบบสองทิศทางระหว่างเว็บเบราว์เซอร์และเซิร์ฟเวอร์ผ่าน WebSocket ใช้สำหรับแจ้งเตือนและแสดงผลข่าวใหม่ทันทีที่ระบบดึงข่าวได้

\subsubsection{KKU Gen.AI}

KKU Gen.AI เป็นบริการโมเดลภาษาขนาดใหญ่ของมหาวิทยาลัยขอนแก่นที่ให้บริการผ่าน API แบบรองรับ OpenAI-compatible ใช้สำหรับการสรุปเนื้อหาข่าวและสกัดประเด็นสำคัญของระบบ

\subsubsection{กลไกแคช (Cache Mechanism)}

กลไกแคชใช้สำหรับจัดเก็บสำเนาข้อมูลข่าวและเนื้อหาที่ดึงมาแล้ว โดยใช้ตัวระบุเฉพาะของข่าว เช่น URL เป็นกุญแจแคช เพื่อหลีกเลี่ยงการดึงข้อมูลซ้ำจากเว็บไซต์ต้นทางในรอบการทำงานถัดไป ช่วยลดภาระของเว็บไซต์ต้นทางและประหยัดทรัพยากรของเซิร์ฟเวอร์

\subsubsection{ระบบคุกกี้ (Cookie System)}

ระบบคุกกี้ใช้สำหรับจัดเก็บข้อมูลสถานะของผู้ใช้และเซสชันการเชื่อมต่อ ทั้งเพื่อคงสถานะการดึงข้อมูลจากเว็บไซต์ที่ต้องใช้คุกกี้หรือการตรวจสอบเบราว์เซอร์ และเพื่อจัดเก็บค่ากำหนดเฉพาะบุคคลและข้อมูลพฤติกรรมการอ่านสำหรับอัลกอริทึมแนะนำข่าว

\subsubsection{อัลกอริทึมแนะนำข่าว (Recommendation Algorithm)}

อัลกอริทึมแนะนำข่าวใช้ข้อมูลความสนใจที่บันทึกไว้ในคุกกี้เพื่อวิเคราะห์และคัดสรรข่าวที่เกี่ยวข้องมากที่สุดมาแสดงเป็นลำดับแรกบนหน้าเว็บ โดยคำนวณคะแนนความสนใจของแต่ละหมวดหมู่จากความถี่การเข้าชม เวลาที่ใช้อ่าน และการกระทำที่แสดงความสนใจ เช่น การกดบันทึกหรือแชร์
